\section{Partitions defined by linear equations}\label{sec:product-tests}

A linear partition records a map \(L:K\to\F^r\) on the ambient key
vector space. Its cells are the fibers of \(L\), intersected with the
Gram reference set. Conditional averages are taken with respect to the
reference probability \(\nu\). Write \(\Pi_L d\) for the function equal
to the mean of \(d\) on each cell.

\subsection{Large Fourier coefficients on a leaf}

Write \(u\) for uniform probability on the ambient key space,
\(a=u(\cK)\), and \(\chi_\xi(x)=(-1)^{\xi\cdot x}\). The Gram count
gives \(a\ge a_0>0\) for all large \(n\). The reference law
\(\nu=u(\,\cdot\mid\cK)\) has exponentially small bias on nonzero characters, uniformly in the
frequency. This follows by expanding the prescribed Gram equations.
For each nonconstant quadratic term the ambient rank gives exponential
cancellation, and the constant term has zero mean against a nonzero
linear character.

Fix a capped leaf. Let \(j=(o,v)\) record the orientation and the query
flags, with probabilities \(\mu_j\). Write \(d_j\) for the conditional key density
relative to \(\nu\), restricted to \(\cK\). Flags of mass zero can be
omitted. A filter assigns weights \(0\le t_j\le\mu_j\) and gives density
\(d_t=\sum_j t_jd_j\). It may be chosen using the opposite leaf.

For any fixed bound on the number of rows, the cell estimate
\eqref{eq:linear-cell-cap} gives constants \(V,B_0\) such that
\begin{equation}\label{eq:linear-density-energy}
  \sum_j\mu_j\|\Pi_Ld_j\|_2^2\le V,\qquad 0\le\Pi_Ld_t\le B_0.
\end{equation}
The constants precede the row bound. Indeed, outside exponentially small
orientation mass, the unfiltered probability of every cell \(C\) is at
most \(B\nu(C)\). If \(a_{v,C}\) is its joint probability with flag
\(v\) and \(r_v\) is the flag probability, then
the contribution of a flag \(v\) with \(r_v>0\) is bounded by
\[
 \sum_C\frac{a_{v,C}^2}{r_v\nu(C)}
 \le B\sum_C\frac{a_{v,C}}{r_v}=B .
\]
Summing over the finitely many flags gives a bound independent of \(L\).
On the exceptional orientations the sum is at most the reciprocal of
the smallest positive cell mass. With a fixed row bound this contributes
\(o(1)\). For \(d_t\) itself, the unfiltered law dominates every filter.
Its density on a cell is at most \(B\) plus the exceptional mass
divided by that cell's reference probability. This gives a common
\(B_0\) once \(n\) is large.

% lax begin Lax793645.LargeFrequencies
\begin{lemma}\label{lem:linear-test}
Fix \(\tau>0\). On a given leaf the number of frequencies \(\xi\) for
which at least one filter has \(\big|\int d_t\chi_\xi\,d\nu\big|>\tau\)
is bounded independently of \(n\).
\end{lemma}
% lax end
% lax begin Lax793645Proofs.Manuscript.LargeFrequencies.filtered_large_spectrum
\begin{proof}
For a finite set \(S\) of frequencies let \(L\) record its characters.
Put \(g_j=\Pi_Ld_j\). Its correlations with those characters are the
same as those of \(d_j\). Parseval applied to \(1_{\cK}g_j\) gives
\[
 a\sum_{\xi\in S}\left|\int d_j\chi_\xi\,d\nu\right|^2
 \le \|g_j\|_2^2 .
\]
If a filter witnesses a correlation greater than \(\tau\), then
Cauchy--Schwarz gives
\(\sum_j\mu_j|\int d_j\chi_\xi\,d\nu|^2>\tau^2\).
Choose an integer \(R>V/(a_0\tau^2)\) and apply
\eqref{eq:linear-density-energy} at that row count. Summing over any
\(R\) distinct such frequencies would exceed \(V\).
\end{proof}
% lax end

% lax begin Lax793645.FilteredQueryPartition
\begin{proposition}\label{prop:all-filter-partition}
For each \(\epsilon>0\), one linear partition of bounded dimension
satisfies \(\|d_t-\Pi_Ld_t\|_1\le\epsilon\) and
\(0\le\Pi_Ld_t\le B_0\) for every filter on the leaf.
\end{proposition}
% lax end
% lax begin Lax793645Proofs.Manuscript.FilteredQueryPartition.one_query_all_filters_linear_partition
\begin{proof}
We choose the partition from the large coefficients of \emph{all}
filtered densities on the leaf. Its size is bounded by the preceding
lemma. To prove it suffices, we show that a bounded test orthogonal
to its characters has small integral against each filtered density.

Put \(\eta=a_0\epsilon/8\). Expand the flags as in
\Cref{thm:leaf-fourier} and divide the error among their fixed number
of terms. Let \(\delta>0\) be the resulting cutoff for an error of
\(\eta\) against any filtered density.
Put \(k=\lceil\delta^{-2}\rceil\) and
\(\tau=a_0\epsilon/(4\cdot2^k)\). Let \(L\) record the set in
Lemma~\ref{lem:linear-test}.

Consider any \(|W|\le1\) orthogonal in \(L^2(\nu)\) to the recorded
characters. Extend \(W\) by zero to \(w=1_{\cK}W\).
By Parseval at most \(k\) coefficients of \(w\) exceed \(\delta\) in
absolute value. Let \(T\) be the conditional average of \(w\) under
\(u\) on the partition they define. Its Fourier coefficients agree with
those of \(w\) on their span and vanish elsewhere. In particular
\(|T|\le1\), its Fourier norm is at most \(2^k\), and its coefficients
vanish at the characters recorded by \(L\). Hence
\(\big|\int T d_t\,d\nu\big|\le2^k\tau\).

Apply \Cref{thm:leaf-fourier} to \(D=1_{\cK}(aW-T)/2\).
This function has modulus at most 1. Its coefficient at \(\xi\) is
\[
  \widehat D(\xi)=\tfrac a2
     \left(\widehat w(\xi)-\int T\chi_\xi\,d\nu\right).
\]
The first term differs from \(\widehat T(\xi)\) by at most \(\delta\).
The second differs from it by at most \(2^k\) times the reference bias,
which is at most \(\delta/2^k\) for large \(n\). Thus every coefficient
of \(D\) has modulus at most \(\delta\). The comparison gives
\[
 a\left|\int W d_t\,d\nu\right|\le 2\eta+2^k\tau.
\]
Finally take \(U=\operatorname{sgn}(d_t-\Pi_Ld_t)\) and
\(W=(U-\Pi_LU)/2\). This test is orthogonal to the recorded characters.
Its integral against \(d_t\) is \(\|d_t-\Pi_Ld_t\|_1/2\), by
self-adjointness of conditional averaging. The chosen constants give the
claim. The pointwise bound follows from \eqref{eq:linear-density-energy}.
\end{proof}
% lax end

\subsection{Clipping follows from approximation}\label{sec:partition-clipping}

Apply Proposition~\ref{prop:all-filter-partition} with all statuses retained. This
gives a bounded approximation \(0\le g\le B_0\) to the unfiltered query
density \(d\). Every flagged query density \(d_f\) satisfies
\(0\le d_f\le d\), even when its flag is not among the recorded statuses.
If \(\|d-g\|_1\le\epsilon\), then
\[
 \int_S d_f\,d\nu\le B_0\nu(S)+\epsilon,
 \qquad
 \int_{\{d_f\ge2B_0\}}d_f\,d\nu\le2\epsilon.
\]
The first inequality follows from \(d\le B_0+|d-g|\). On the set
in the second inequality, \(d_f\le2(d-g)\le2|d-g|\).
We can use the same cutoff \(2B_0\) for every desired tail error.

\subsection{A common set of linear coordinates}

Work first with one tuple format. For each leaf \(\lambda\), collect the
frequencies of its partitions for the options of that format. Their span
\(H_\lambda\) has dimension at most a fixed \(r\). Repeat the argument
for the finitely many formats, dividing the error budget among them.

% lax begin Lax793645.SharedFrequencies
\begin{lemma}\label{lem:common-frequencies}
For every \(\eta>0\), there is a space \(H\) of dimension at most
\(r2^r/\eta\), rounded up, such that two independent leaves satisfy
\begin{equation}\label{eq:frequency-intersection}
 H_\lambda\cap H_{\lambda'}\subseteq H
\end{equation}
except with probability at most \(\eta\).
\end{lemma}
% lax end
% lax begin Lax793645Proofs.Manuscript.SharedFrequencies.shared_leaf_frequencies
\begin{proof}
Call a frequency frequent if it belongs to \(H_\lambda\) with
probability greater than \(\eta/2^r\). Let \(H\) be their span and
choose a basis from the frequent frequencies. At most \(r\) of these basis vectors belong to any \(H_\lambda\),
because they are linearly independent. The expected number it
contains is at least \((\dim H)\eta/2^r\), giving the dimension bound.
Fix the first leaf. Each of its at most \(2^r\) frequencies outside
\(H\) belongs to the independent second leaf with probability at most
\(\eta/2^r\). Sum these probabilities.
\end{proof}
% lax end

Keep the two bounded approximants \(g_\lambda,g_{\lambda'}\) on their
original leaf partitions. Their Fourier coefficients are supported on
\(H_\lambda,H_{\lambda'}\), respectively. Under uniform ambient
probability only common frequencies contribute to their inner product.
On pairs satisfying \eqref{eq:frequency-intersection}, this gives
\begin{equation}\label{eq:common-projection-product}
 \int g_\lambda g_{\lambda'}=
 \int(\Pi_Hg_\lambda)(\Pi_Hg_{\lambda'}).
\end{equation}
For the Gram reference law the error is \(o(1)\), uniformly for bounded
functions on the two leaf partitions. Expand them in their bounded
number of frequencies and use the exponentially small bias of each
nonzero character. The common cell weights obey the same estimate.
The cell bound at the common row count bounds the actual projected
densities. Conditional averaging contracts their \(L^1\) errors.
Thus, with errors chosen small enough, every permitted pair of filters has
bounded approximants \(g_A,g_B\). Put \(\widetilde d_A=\min(d_A,g_A)\)
and similarly on the other side. Then
\begin{equation}\label{eq:bounded-overlap-transfer}
 \int d_Ad_B\ge\int\widetilde d_A\widetilde d_B
 \ge\int g_Ag_B-B_0(\|d_A-g_A\|_1+\|d_B-g_B\|_1).
\end{equation}
Together with \eqref{eq:common-projection-product}, this bounds the average
overlap of the common projections by the actual average overlap plus an
arbitrarily small error. Exceptional leaf pairs cost at most \(B_0^2\eta\).
All partitions were chosen before the filters, so filters depending on
both leaves are allowed. Independence is used only for the two original
leaf draws in Lemma~\ref{lem:common-frequencies}.

\subsection{Testing a tuple one position at a time}

Let \(L\) record the common coordinates on a tuple \(k=(k_i)_i\).
Record \(L(\iota_i k_i)\) in the unary partition at position \(i\), where
\(\iota_i\) inserts one position and puts zero elsewhere. Then
\begin{equation}\label{eq:exact-product}
 \varphi(Lk)=\sum_{(a_i)}\varphi\Bigl(\sum_i a_i\Bigr)
                 \prod_i1_{\{L(\iota_i k_i)=a_i\}}.
\end{equation}
For a given tuple exactly one term survives, the one with
\(a_i=L(\iota_i k_i)\) at every position. The number of terms depends
only on the number of rows and positions. These unary partitions
can be combined across all templates and labels.

For a product of unary tests, expand the finite binary prescriptions.
The constant character contributes their fixed binary probability
\(\kappa>0\) times the product of the unary accepted measures.
Every nonconstant character has a nominal bilinear block of rank at least
\(n/5-O(1)\), by the explicit mixers. The unary statuses have bounded
rank and affect only the \(O(1)\) term. Cross pairings between different
endpoint frames satisfy the analogous rank bound outside exponentially
small mixed mass. Under the raw law, failure requires
\(\Omega(n^2)\) independent bit conditions. Its probability remains
exponentially small after multiplication by the unit density bound. Apply the Walsh bound and sum the
bounded number of characters. This gives product comparison on all the
tests in \eqref{eq:exact-product}.

For an individual cell we also need concentration under the prepared
mixed law. The second moment of its query probability is the probability
that two independent nominal inputs at the \emph{same} frame both give
a hit. For independent nominal columns the image law is uniform among independent ambient lists with the
prescribed Gram data. In the expansion of the joint test, a character
that uses a pairing between the two inputs has a nondegenerate ambient
block. Its contribution to the variance is exponentially small. A
dependence of the nominal columns already has exponentially small
probability by the free \(Z\) coordinates. For any fixed family of
cell tests and any fixed accuracy this gives an exceptional probability
of at most \(2^{-cn}\) under the prepared marginal, for some \(c>0\).

To prescribe basis values as well, use the change of \(Z\) coordinates
from \Cref{sec:moments}. Outside the exceptional frames the bounded
space of basis forms retains the rank needed for unary mixing in the
cell tests. Fourier inversion then gives a fixed positive proportion
for each allowed basis assignment on the cells of sufficient mass.
These concentration estimates use the prepared marginal bound
\(2^{n/32}\mu\). On a single normalized leaf we use
Theorem~\ref{thm:leaf-fourier} and Proposition~\ref{prop:all-filter-partition}
to obtain the simultaneous approximation for its filters.
